% shared by main.tex (IEEEtran) and access/main.tex (IEEE Access). Numbers come from generated/ only.
\section{Introduction}
Language models that write code~\cite{chen2021codex,li2022alphacode,austin2021mbpp} have billions of parameters and run in data centres. Many IoT settings cannot rely on that: field devices are often offline, data may not leave the device, and the hardware has megabytes, not gigabytes, of memory~\cite{warden2019tinyml,david2021tflm,lin2020mcunet}. Even ``small'' language models such as TinyLlama~\cite{zhang2024tinyllama} and MobileLLM~\cite{liu2024mobilellm} have hundreds of millions of parameters. At the other extreme, TinyStories~\cite{eldan2023tinystories} showed that models with a few million parameters can produce fluent text in a narrow domain. This raises a practical question for embedded systems: \emph{what can a code model that fits in microcontroller flash actually do, and which inexpensive techniques make it more reliable?}

Small models make many syntax errors, and a compiler can detect them for free. Feeding compiler diagnostics back to a model is an established idea, from DeepFix~\cite{gupta2017deepfix}, DrRepair~\cite{yasunaga2020drrepair} and Break-It-Fix-It~\cite{yasunaga2021bifi} to self-debugging with large models~\cite{chen2024selfdebug,madaan2023selfrefine,shinn2023reflexion}. Whether it pays off is less clear: Olausson et al.~\cite{olausson2024selfrepair} found that, at an equal sampling budget, self-repair is often no better than drawing more samples. We test this question at a scale four to five orders of magnitude smaller than prior work, where every extra generation would cost real time on a battery-powered device.

\textbf{Contributions.}
\begin{enumerate}
\item \textbf{A leak-free embedded code benchmark.} \nfamtrain{} training and \nfamheld{} held-out task families (sensor scaling, filtering, checksums, bit manipulation, lists, strings and recursion), each with a reference solution and executable tests. A script proves that no test instruction appears in training and that held-out families are never trained on (Section~\ref{sec:bench}).
\item \textbf{Microcontroller-sized models with an exact budget.} Transformers of 0.53 and 1.05\,M parameters (and a 3.05\,M configuration) with a 2{,}048-token tokenizer; every parameter is accounted for in Table~\ref{tab:params}. The 0.53\,M model needs \dsmflashtotalbytes{} bytes of flash including the tokenizer.
\item \textbf{A controlled comparison of cheap reliability techniques} under an equal budget of four generations: signature anchoring, REF with the full message, a line-only message or a constant placeholder, and resample-until-compiles, plus training-data ablations, two model sizes with three seeds each.
\item \textbf{An end-to-end deployment path.} A dependency-free INT8 C engine with a C tokenizer, checked against PyTorch on the host; ESP-IDF firmware that builds for the ESP32-S3 and runs in a host simulation; and, not yet built or run on a device, a MicroPython module, a REF loop that uses MicroPython's own \texttt{compile()} as the gate, and an offline ``field assistant'' script.
\item \textbf{Negative results reported as such.} The model does not generalise to unseen task families or to MBPP, and replacing the error message with a constant placeholder does not measurably change repair accuracy.
\end{enumerate}

\section{Related Work}
\textbf{Small and on-device language models.} TinyML deploys neural networks on microcontrollers through quantisation~\cite{jacob2018quant}, specialised runtimes~\cite{david2021tflm} and joint architecture--runtime design~\cite{lin2020mcunet}, with standard benchmarks for vision and audio~\cite{banbury2021mlperftiny}. For language, TinyStories~\cite{eldan2023tinystories} and phi-1~\cite{gunasekar2023textbooks} show that narrow, high-quality data lets small models punch above their size; TinyLlama~\cite{zhang2024tinyllama} and MobileLLM~\cite{liu2024mobilellm} target phones rather than microcontrollers. We are not aware of code-generation results for models small enough to run from microcontroller flash.

\textbf{Code generation.} Large code models are evaluated by functional correctness on held-out problems such as HumanEval~\cite{chen2021codex} and MBPP~\cite{austin2021mbpp}, often with sample-then-filter decoding~\cite{li2022alphacode}. Smaller pretrained encoder--decoders such as CodeT5~\cite{wang2021codet5} still have tens of millions of parameters.

\textbf{Repair from compiler or execution feedback.} DeepFix~\cite{gupta2017deepfix} and DrRepair~\cite{yasunaga2020drrepair} train models to fix compiler errors; Break-It-Fix-It~\cite{yasunaga2021bifi} learns to repair real broken code with a learned breaker and verified fixes, a related idea; Section~\ref{sec:size} reports why a self-generated variant was not feasible here; Self-Debug~\cite{chen2024selfdebug}, Self-Refine~\cite{madaan2023selfrefine} and Reflexion~\cite{shinn2023reflexion} prompt large models with feedback. Olausson et al.~\cite{olausson2024selfrepair} show that repair must be compared with resampling at equal cost; we follow that protocol. Automated program repair more broadly is surveyed in~\cite{legoues2019apr}.

\section{Benchmark}\label{sec:bench}
\textbf{Task families.} A family is a parameterised task: four hand-written natural-language phrasings, a canonical reference solution, and a generator of test inputs. A concrete task samples a function name and argument names from small pools and the family's constants (e.g.\ ADC resolution and reference voltage, filter coefficient, threshold, bit index). Expected outputs are always produced by executing the reference solution. There are 64 families in five categories: \emph{sensor/IoT} (ADC scaling, unit conversion, clamping, range mapping, moving and exponential averages, RMS, median, peak counting, debounce, hysteresis, XOR/additive checksums, byte packing, bit tests), \emph{math}, \emph{list}, \emph{string} and \emph{recursion}. All reference solutions are written to be MicroPython-compatible (we have not executed them under MicroPython).

\textbf{Splits.} With a fixed seed and stratified by category, \nfamheld{} families (25\,\%) are held out and never used for training. For the remaining \nfamtrain{} families, phrasing~4 is never used in training. \emph{In-family} test (\ntestid{} tasks): training families, phrasing~4 only, freshly sampled names and constants. \emph{Out-of-family} test (\ntestood{} tasks): held-out families, any phrasing. Model selection uses a separate \emph{development} split (Section~\ref{sec:select}). We also report MBPP (500 test problems)~\cite{austin2021mbpp} as a fully out-of-domain check.

\textbf{Leakage audit.} \texttt{scripts/check\_leakage.py} verifies: zero exact or normalised instruction overlap between training and both test splits; zero test tasks from a family seen in training for the out-of-family split; zero in-family test instructions sharing any 8-gram with training instructions (\ngramood{} out-of-family tasks, all from \texttt{list\_max}, share an 8-gram such as ``of the non-empty list items using a loop'' with the trained \texttt{list\_min} phrasing). By construction, \codeoverlapid{} of the \ntestid{} in-family reference solutions also occur in training: the in-family split tests whether the model recognises a known task from an unseen phrasing, and the out-of-family split tests whether it can write code it has never seen.

\textbf{Metric.} A task passes if the generated program compiles and returns the expected value (floats within $10^{-6}$) on all five test inputs, executed in a sandboxed child process with a 2\,s time limit. We report mean\,$\pm$\,s.d.\ over three training seeds; per-run 95\,\% Wilson intervals are in the released result files.

\section{Model and Training}
\textbf{Architecture.} A decoder-only Transformer~\cite{vaswani2017attention} with RMSNorm~\cite{zhang2019rmsnorm}, rotary position embeddings~\cite{su2024roformer}, a GELU MLP, no biases, and input/output embeddings tied~\cite{press2017tying}. The architecture is deliberately plain so that it can be re-implemented in about 300 lines of C. Table~\ref{tab:params} gives the exact parameter budget. With a small vocabulary the embedding is about a quarter of the 1.05\,M model (37\,\% of the 0.53\,M model), which is what makes a sub-megabyte INT8 image possible.

\begin{table}[t]\centering\caption{Model sizes (vocabulary 2{,}048, context 256 tokens, MLP ratio 4). Every parameter is counted.}\label{tab:params}
\setlength{\tabcolsep}{3.5pt}
\begin{tabular}{lrrrrrrr}\toprule
 & $L$ & $d$ & $H$ & Embedding & Blocks & Total & MiB (INT8)\\\midrule
\inputtab{generated/tab_params}
\bottomrule\end{tabular}\end{table}

\textbf{Tokenizer.} A byte-level BPE~\cite{sennrich2016bpe} with 2{,}048 entries trained on the training corpus only. Because it is byte-level, any input can be encoded; unseen words simply fall back to short pieces.

\textbf{Training data.} (i) \emph{Generation pairs} (\ntraingen{}): instruction $\rightarrow$ code. (ii) \emph{Synthetic repair triples} (\ntrainrepair{}): instruction, a corrupted program, the exact CPython \texttt{SyntaxError} message, and the correct program. Corruptions delete a colon or bracket, break indentation, misspell a keyword, drop an operator or duplicate a token, and are kept only if they really fail to compile. We also tried training-only instruction augmentation (random function/argument names, varied request openings such as ``Make a function'' or ``I need a function''); it lowered accuracy on the development split, so the final recipe does not use it (Section~\ref{sec:select}). The model is trained with the loss on the answer tokens only, on an equal mix of generation and repair examples, with AdamW, a cosine schedule, dropout 0.1, batch 48 and 4{,}000 steps (3{,}000 for the 3.05\,M model, 2{,}500 for development-split runs) on a CPU without a GPU (17--33\,min per model).

\section{Inference Strategies}
All strategies use at most four generations per task, so they are compared at an equal worst-case number of generations (repair prompts are longer, so compute is not identical).

\textbf{Signature anchoring (SA).} Every request names the function to write. SA extracts the first \texttt{name(args)} from the request and forces the answer to begin with \texttt{def name(args):}; the model only completes the body. SA adds no parameters and saves the tokens of the header.

\textbf{Recursive Error-Feedback (REF).} Algorithm~\ref{alg:ref}: if the draft does not compile, the model is prompted with the instruction, the faulty draft and the error message, using the same special-token format as the repair training data. If greedy repair returns the draft unchanged, one sample at temperature 0.7 is drawn instead. We also test a second gate, a static undefined-name check described in Section~\ref{sec:fail}, applied after \texttt{compile()}. To isolate what the message contributes we compare three message variants: the full CPython message; a line-only message (``\texttt{SyntaxError: invalid syntax (line $n$)}''), which approximates what MicroPython reports (not checked on MicroPython); and a constant placeholder (``\texttt{SyntaxError}'') instead of the message.

\textbf{Resample-until-compiles.} The equal-budget baseline of~\cite{olausson2024selfrepair}: while the program does not compile, draw a fresh sample (temperature 0.7) from the original prompt.

\begin{algorithm}[t]\caption{Signature anchoring + REF}\label{alg:ref}
\begin{algorithmic}[1]
\Require instruction $I$, budget $N=3$ repairs
\State $p \gets$ \texttt{def name(args):} parsed from $I$ \Comment{SA}
\State $C \gets p \,\|\, \textsc{Greedy}(\langle q\rangle I \langle c\rangle p)$
\For{$n = 1..N$}
  \State $E \gets \textsc{Gate}(C)$ \Comment{\texttt{compile()}, optionally + name check}
  \State \textbf{if} $E=\emptyset$ \textbf{return} $C$
  \State $P \gets \langle q\rangle I \langle b\rangle C \langle r\rangle E \langle c\rangle p$
  \State $C' \gets p \,\|\, \textsc{Greedy}(P)$; \textbf{if} $C'=C$: $C' \gets p \,\|\, \textsc{Sample}_{0.7}(P)$
  \State $C \gets C'$
\EndFor
\State \Return $C$ if $\textsc{Gate}(C)=\emptyset$, else \textsc{Failure}
\end{algorithmic}\end{algorithm}

\section{Deployment Path for the ESP32-S3}
\textbf{INT8 engine.} Weights are quantised per output row to symmetric INT8 with FP32 scales; activations and the KV cache stay FP32 (the ESP32-S3 has a single-precision FPU). The engine is portable C99, performs no allocation in the decode loop and reads weights in place from a constant blob, so on the ESP32-S3 they are designed to stay memory-mapped in flash (not yet verified on hardware). The tokenizer (GPT-2 pre-tokenisation plus rank-ordered BPE merges) is implemented in the same file. On the host, the C tokenizer reproduces the reference tokenizer on all 18{,}524 instruction and code strings of the benchmark, the generation part of the training corpus and MBPP, and on one test prompt the C logits of the 1.05\,M model match PyTorch with identically quantised weights to within $1.5\times10^{-5}$. A JavaScript port of the engine for a browser playground reproduces the host simulation's output on all \ntestid{}+\ntestood{} test prompts (0.53\,M model; \texttt{scripts/check\_js\_engine.sh}).

\textbf{Memory.} Flash holds the model (\depintmodelbytes{} bytes) and tokenizer (\deptokenizerbytes{} bytes). RAM holds the runtime state of \depstatebytes{} bytes, of which \depkvcachebytes{} bytes are the FP32 KV cache for the full 256-token context; this requires PSRAM (e.g.\ an 8\,MB-PSRAM module); Table~\ref{tab:deploy} gives the same figures for the 0.53\,M model, which is the one embedded in the firmware. One generated token costs \depmacsfirst{} multiply-accumulates at the first position and \depmacslast{} at position 256; the projection onto the 2{,}048-token vocabulary accounts for 262{,}144 of these (about one fifth to one quarter), and attention over a full 256-token context for another fifth.

\textbf{Firmware, simulation and application.} The ESP-IDF firmware embeds the model and the benchmark prompts and is written to run SA + greedy decoding on every prompt and print the code and timings over UART, where a script re-scores the log with the paper's unit tests; it has been built in CI but not yet run on a board. The same \texttt{app\_main()} also runs on a PC with emulated ESP32-S3 heap budgets (300\,KiB internal, 8\,MiB PSRAM; host simulation); a Wokwi configuration is included but has not been run. Built in CI with ESP-IDF v5.4.2 for the ESP32-S3 with the 0.53\,M model embedded, the firmware image is \idfappbytes{} bytes (\idfapppct\,\% of its \idfapppartitionbytes{}-byte app partition): \idfflashdatabytes{} bytes of read-only data in flash, most of it the model and tokenizer, and \idfflashcodebytes{} bytes of code; static use of internal RAM is \idfdiramusedbytes{} of \idfdiramtotalbytes{} bytes (\idfdirampct\,\%). For REF on the device we provide a MicroPython user C module and a REF loop in plain MicroPython that uses MicroPython's own \texttt{compile()} as the gate; neither has yet been built or run on a device. On top of it, a \emph{field assistant} script (untested on hardware) is intended to let a technician describe a sensor function over the serial console and have the device generate, gate and install it, with a helper that applies it to an ADC channel; it uses a separate application model trained on all families and phrasings, which is not evaluated in this paper.

\section{Results}
\subsection{Model selection on the development split}\label{sec:select}
Our first two attempts were evaluated directly on the test split (Table~\ref{tab:abl}, first two rows; re-scored on the final split). The first used 500 draws per family, 3{,}000 steps and no dropout (its training corpus was not retained); the second used varied request openings and replaced each function and argument name by a random identifier with probability 0.5, with dropout 0.1 and 4{,}000 steps, so the two differ in several ways. The second was clearly worse with signature anchoring, so for the augmentation decision we built a separate \emph{development} split: models are trained on phrasings 1--2 only and evaluated on \ndev{} tasks written with phrasing~3 of the training families. Test phrasing~4 is never seen. In this single-seed comparison (Table~\ref{tab:dev}) augmentation hurts. Varying how a request opens lowers SA+REF accuracy from \devNoaugarefpass\,\% to \devOpenersarefpass\,\%, and random identifiers lower it further, to \devNamesfiftyarefpass\,\% at 50\,\%. We conjecture that descriptive names such as \texttt{celsius\_to\_fahrenheit} are a strong cue to the intended task, and that randomising them removes it. The final recipe therefore uses no instruction augmentation. All results below use this recipe; the augmentation choice was made on the development split, but the undefined-name gate (Section~\ref{sec:fail}) and the choice of deployed model were informed by test-split results.

\begin{table}[t]\centering\caption{Development-split selection (\ndev{} tasks, pass \%, seed 0, 2{,}500 steps).}\label{tab:dev}
\begin{tabular}{lrrr}\toprule
Training data & Greedy & SA & SA+REF\\\midrule
\inputtab{generated/tab_dev}
\bottomrule\end{tabular}\end{table}

\subsection{Main results}
Table~\ref{tab:main} reports the 1.05\,M model over three seeds; the 0.53\,M model is compared in Section~\ref{sec:size}.

\begin{table*}[t]\centering
\caption{Pass rate and syntactic validity (\%) of the 1.05\,M-parameter model, mean\,$\pm$\,s.d.\ over \nseeds{} seeds. ``Gens'' is the mean number of generations used (budget 4).}\label{tab:main}
\begin{tabular}{lccccc}\toprule
 & \multicolumn{2}{c}{In-family (\ntestid{} tasks)} & \multicolumn{2}{c}{Out-of-family (\ntestood{} tasks)} & \\
Strategy & Valid & Pass & Valid & Pass & Gens\\\midrule
\inputtab{generated/tab_main}
\bottomrule\end{tabular}\end{table*}

\textbf{Signature anchoring is the largest single gain.} Greedy decoding passes \testidgreedypass\,\% of in-family tasks; forcing the header raises this to \testidanchorpass\,\% (seed range \testidanchorpassmin--\testidanchorpassmax\,\%) at zero cost. Without SA the model often writes a correct body under a slightly different name or argument list, which fails every test.

\textbf{Compiler feedback mostly fixes syntax, not behaviour.} REF raises syntactic validity on the in-family split from \testidanchorvalid\,\% to \testidarefvalid\,\%, and on the out-of-family split from \testoodanchorvalid\,\% to \testoodarefvalid\,\%. The pass rate barely moves (\testidanchorpass\,\% $\rightarrow$ \testidarefpass\,\%), because most failures are semantic: a wrong constant, an undefined name, a wrong loop bound (Section~\ref{sec:fail}).

\textbf{Replacing the error message changes little, and repair is not better than resampling.} With the full CPython message, a line-only message and a constant placeholder, SA+REF reaches \testidarefpass, \testidareflinepass{} and \testidarefnoerrpass\,\%. Resample-until-compiles reaches \testidaresamplepass\,\%. All four differences are far smaller than the seed-to-seed spread (s.d.\ \testidarefpasssd{} points). This is consistent with the finding of~\cite{olausson2024selfrepair} that at equal cost self-repair gains little over resampling, here at a scale about $10^{4}$ times smaller. In our setting, replacing the message with a placeholder made no measurable difference, whereas~\cite{olausson2024selfrepair} found feedback quality to be a bottleneck for large models.

\textbf{A second, equally cheap gate helps a little.} An undefined-name check (Section~\ref{sec:fail}) catches failures that \texttt{compile()} accepts. Gating REF on \texttt{compile()} plus this check gives the best in-family result, \testidareflintpass\,\%, using \testidareflintgens{} generations on average. The gain over SA alone (\testidanchorpass\,\%) is consistent across the three seeds but smaller than the seed-to-seed spread (s.d.\ \testidareflintpasssd{} points).

\textbf{Out-of-family generalisation is essentially absent.} On the \ntestood{} tasks from held-out families every strategy passes between \testoodgreedypass{} and \testoodaresamplelintpass\,\% (mean over seeds). REF still makes most of these programs \emph{compile} (\testoodrefvalid\,\% valid without SA). The model writes fluent Python, but for tasks it was not trained on, that Python does the wrong thing.

\subsection{Ablations}
\begin{table*}[t]\centering\caption{Ablations (pass \%, one run each). ID: in-family split; OOD: out-of-family split. SA+REF+name: REF gated by \texttt{compile()} and the undefined-name check.}\label{tab:abl}
\begin{tabular}{lrrrrrr}\toprule
 & & \multicolumn{4}{c}{In-family (ID)} & OOD\\\cmidrule(lr){3-6}
Variant & Params & Greedy & SA & SA+REF & SA+REF+name & SA+REF+name\\\midrule
\inputtab{generated/tab_ablation}
\bottomrule\end{tabular}\end{table*}

\textbf{Repair data may also help generation} (Table~\ref{tab:abl}).\label{sec:size} In a single run, removing the repair examples lowers even plain greedy accuracy (\ablGenonlygreedypass\,\% vs.\ \ablFinalgreedypass\,\%, seed 0), which suggests that repair triples act as extra supervision; with one seed and a greedy seed s.d.\ of \testidgreedypasssd{} points this is not conclusive. Adding undefined-name repair examples does not help: \ablNamerepareflintpass\,\% with the two-gate REF, against \ablFinalareflintpass\,\% for the final recipe (one seed each, a difference within the seed spread).

\textbf{Learning from the model's own errors does not apply here.} We also tried the Break-It-Fix-It idea~\cite{yasunaga2021bifi} of learning from the model's own mistakes, but on training prompts the model produced only \selfrepairpairs{} syntax error in \selfrepairsamples{} samples at temperature 1.0: its errors appear only on unseen wordings, which we cannot use for training without leaking the test phrasing.

\textbf{Half a megabyte is enough.} Model size matters less than one might expect. Over \seedsSmalln{} seeds, the 0.53\,M model reaches \seedsSmalltestidanchor\,$\pm$\,\seedsSmalltestidanchorsd\,\% with SA and \seedsSmalltestidareflint\,$\pm$\,\seedsSmalltestidareflintsd\,\% with SA+REF and both gates, against \seedsMediumtestidanchor\,$\pm$\,\seedsMediumtestidanchorsd{} and \seedsMediumtestidareflint\,$\pm$\,\seedsMediumtestidareflintsd\,\% for the 1.05\,M model\ifdefined\ablLargeanchorpass; the 3.05\,M model, trained for 3{,}000 steps, reaches \ablLargeanchorpass\,\% with SA (one seed)\fi. For a closed task catalogue, half a megabyte of INT8 weights is enough. Apart from the size comparison, the ablation rows are single runs, and the seed spread of the final recipe (s.d.\ \testidanchorpasssd{} points) is larger than most of these differences, so we do not claim any of them as significant.

\textbf{Quantisation is nearly lossless (seed 0).} Running the full benchmark through the INT8 C engine instead of FP32 PyTorch gives \ablIntanchorpass\,\% with SA and \ablIntareflintpass\,\% with SA+REF and both gates for the 1.05\,M model (FP32: \ablFinalanchorpass{} and \ablFinalareflintpass\,\%), and \ablSmallintanchorpass{} and \ablSmallintareflintpass\,\% for the 0.53\,M model (FP32: \ablSmallanchorpass{} and \ablSmallareflintpass\,\%).

\subsection{Failure analysis}\label{sec:fail}
For seed~0 with SA+REF on the in-family split, the \failidtotal{} failures split into \failidruntime{} runtime exceptions, \failidwrong{} wrong values, \failidpartly{} partly correct, and \failidsyntax{} programs that still do not compile. Of the runtime exceptions, \failidnameerr{} are \texttt{NameError}s: the model reuses a variable name from a related family (for instance returning \texttt{s} where a constant upper bound was required). This observation motivated the undefined-name gate: a 118-line static check, written without \texttt{re} or \texttt{ast} so that it is intended to run in MicroPython (not yet run there), which flags any identifier that is read but never bound. It produces no false positives on all 5{,}581 reference solutions and flags every \texttt{NameError} produced by the model.

\subsection{Out-of-domain: MBPP}
On the 500 MBPP test problems (prompt plus ``The function is called \textit{name}.''), the 1.05\,M model (seed 0 only) writes syntactically valid Python for \mbppgreedyvalid\,\% of problems (\mbpparefvalid\,\% with SA+REF) and passes \mbppgreedypassed{} (greedy) and \mbpparefpassed{} (SA+REF) of them. This is the expected outcome for a model trained only on \nfamtrain{} templated task families, and it marks the boundary of what such a model should be used for.

\subsection{Deployment measurements (host simulation)}
The firmware path is checked end to end on the host: running the unmodified \texttt{app\_main()} with the INT8 blob on all \ntestid{}+\ntestood{} prompts gives a SA-greedy pass rate of \dephostsimgreedypasspct\,\%, identical to the FP32 PyTorch result on the same prompts. In the host simulation, the state allocation succeeds with 8\,MiB of emulated PSRAM and fails with none, as the memory budget predicts. On one x86 core the engine decodes at \dephostdecodemspertoken\,ms per token. Because the 0.53\,M model is as accurate as the 1.05\,M model (Section~\ref{sec:size}) and needs half the flash and compute, it is the one we select for deployment. \ifdefined\devdecode
On an ESP32-S3 at 240\,MHz (board run, \texttt{esp32/run\_lab.sh board}), the 0.53\,M model decodes at \devdecode\,ms per token and completes a median task in \devendtoend\,s. Every generated program is printed over UART and re-scored with the same unit tests: \devpass\,\% of all \devn{} prompts pass, so the on-device accuracy is measured, not assumed.
\else
We have not yet measured the model on ESP32-S3 hardware. For orientation only, a simple estimate assuming 4 cycles per INT8 multiply--accumulate at 240\,MHz, plus streaming all weights from external flash at 40\,MB/s for every token, gives, at mid-context (position 128), \estlo\,ms per token for compute alone and \esthi\,ms if all weights are also streamed from flash, for the 0.53\,M model; a board run (\texttt{esp32/run\_lab.sh board}) replaces this estimate with measured latency and on-device accuracy.
\fi

\begin{table}[t]\centering\small\caption{Deployment. Host values are measured with the same C engine and firmware code; ESP32-S3 values \ifdefined\devdecode come from the board run\else are an analytical estimate (see text)\fi{} for the 0.53\,M model.}\label{tab:deploy}
\setlength{\tabcolsep}{3pt}
\begin{tabular}{lrr}\toprule
 & 1.05\,M (M) & 0.53\,M (S)\\\midrule
Flash: model + tokenizer (B) & \depflashtotalbytes & \dsmflashtotalbytes\\
RAM: state incl.\ KV cache (B) & \depstatebytes & \dsmstatebytes\\
MACs/token, position 1 & \depmacsfirst & \dsmmacsfirst\\
MACs/token, position 256 & \depmacslast & \dsmmacslast\\
Host decode (ms/token) & \dephostdecodemspertoken & \dsmhostdecodemspertoken\\
Host end-to-end, median (ms) & \dephostendtoendmsmedian & \dsmhostendtoendmsmedian\\
Host SA-greedy pass, all 624 (\%) & \dephostsimgreedypasspct & \dsmhostsimgreedypasspct\\\midrule
\ifdefined\devdecode
ESP32-S3 decode (ms/token) & & \devdecode\\
ESP32-S3 prefill (ms/token) & & \devprefill\\
ESP32-S3 end-to-end, median (s) & & \devendtoend\\
ESP32-S3 SA-greedy pass, all 624 (\%) & & \devpass\\
\else
ESP32-S3 decode, \emph{estimate} (ms/token) & & \estlo--\esthi\\
ESP32-S3 on-device accuracy & & not measured\\
\fi
\bottomrule\end{tabular}\end{table}

\section{Discussion and Limitations}
\textbf{What a million parameters can do.} The model is a reliable \emph{recall-and-adapt} engine. Given an unseen wording of a familiar task, it fills in the right constants, names and structure for roughly four out of five requests, and two free gates make almost all of its output compile. It does not compose new algorithms. This makes it a good fit for a closed catalogue of device-side utilities (unit conversions, filters, thresholds, checksums, register bit operations) that a technician adapts in the field. It is not a replacement for a general coding assistant.

\textbf{Cheap techniques first.} At this scale the decoding procedure matters more than the error feedback: SA gives the largest gain, and replacing the compiler message with a placeholder changed little. In our runs, almost all repairs that succeeded did so on the first gated retry, so SA plus a single retry may suffice on a device with a strict latency budget.

\textbf{Limitations.} (i) The benchmark is synthetic and templated. In-family tests change the wording, names and constants but reuse known solution templates (\codeoverlapid{} of \ntestid{} reference solutions occur in training). (ii) Each family has only four phrasings, and the model is sensitive to wording; real user requests will be more varied than our test phrasing. (iii) Seed-to-seed variance is large (in-family SA pass between \testidanchorpassmin{} and \testidanchorpassmax\,\%), so single-seed ablations are indicative only. (iv) Our first two training attempts were evaluated on the test split before we introduced the development split; we report them rather than hide them. (v) The \texttt{compile()} and name gates check form, not function; generated code must still be tested before it controls hardware. (vi) \ifdefined\devdecode The ESP32-S3 numbers come from a single board run.\else Nothing has been run on ESP32-S3 hardware or under MicroPython: on-device latency is an estimate, on-device accuracy is unmeasured, and the MicroPython module and field assistant are untested.\fi

\section{Conclusion}
We built an openly released pipeline for code generation with Transformers of 0.53 and 1.05\,M parameters. The 0.53\,M model fits in about half a megabyte of microcontroller flash; on unseen wordings of known embedded tasks it passes the unit tests on \seedsSmalltestidanchor\,\% of tasks with signature anchoring and \seedsSmalltestidareflint\,\% with gated repair (1.05\,M model: \testidanchorpass{} and \testidareflintpass\,\%). It does not generalise to new task families. Compiler feedback makes its output compile but adds little correctness, and replacing the message with a placeholder changed little. We hope the benchmark, the leakage audit, the INT8 engine and the negative results help others measure on-device code models honestly.

\section*{Reproducibility}
All numbers are generated by \texttt{scripts/make\_tables.py} from the result files written by the evaluation scripts; \texttt{scripts/reproduce.sh} reruns data generation, the leakage audit, the main training and evaluation runs and the host simulation; the remaining runs (extra seeds of the 0.53\,M model, ablations, MBPP, INT8 evaluation and parity tests) are in the other scripts under \texttt{scripts/}, and the training corpora of the first two attempts were not retained.

